\section{Ontogeny Recapitulates Phylogeny}

\begin{quotationx}
Judgment is between Apprehension and Appetite. First, a man apprehends a thing; then he judges it to be good or bad; then he pursues or avoids it accordingly. If therefore the Judgment guides the Appetite wholly, and in no way is forestalled by the Appetite, then is the Judgment free. But if the Appetite in any way at all forestalls the Judgment and guides it, then the Judgment cannot be free; it is not its own: it is captive to another power. Therefore the \textbf{brute beasts} cannot have freedom of Judgment; for in them the Appetite always forestalls the Judgments.

It is therefore again manifest that this liberty, or this principle of all our liberty, is the greatest gift bestowed by God on mankind: by it alone we gain happiness as \textbf{men}: by it alone we gain happiness elsewhere as \textbf{gods}.

\flright{\textsc{Dante Alighieri}, \emph{De Monarchia}}
\end{quotationx}

In the 19\textsuperscript{th} century, \textbf{Ernst Haeckel} popularized ``\textbf{law or recapitulation}'' by giving it an alleged scientific foundation. This law states that ontogeny recapitulates phylogeny, or in other words, the development of an individual organism follows the stages of the evolution of the species. To buttress his case, Haeckel provided tendentious drawings of embryonic development that still show up in biology textbooks to this day, despite being discredited by biologists. Nevertheless, this law has had wide influence, not least of all on \textbf{Rudolf Steiner} who referred to it with approval.

As understood in the strictly material sense, the law as stated is false. However, there is a metaphysical version of the law that is more promising and was actually held by \textbf{Thomas Aquinas}. In this exposition, the ontogeny of the individual follows the development of the soul from the human to the plant to the animal and to the human kingdoms. The embryo begins as a biological and material being. It then develops the vegetative soul which governs its growth, nourishment, and development. Next comes the sensitive soul as the embryo feels its environment and initiates motion. Finally, in the human case, there is the rational soul which, as the form of the body, unifies the other souls and body. This is not a natural process. Since a material process cannot cause an immaterial substance as Aquinas claims.

This, however, is not the same as evolution. Priority does not just refer to temporal sequence. As long as one whose thinking is limited to time, then he cannot think like a metaphysician as \textbf{Rene Guenon} asserts. A more fundamental notion is ontological priority. Hence, we can say that the idea of the rational soul precedes the sensitive and vegetative souls and the body, although they appear in time in the reverse sequence. Here Aquinas' explanation seems somewhat shaky. He claims that the rational soul is created at the time it is infused with the body, replacing and including the other souls. First of all, phenomenological analysis shows that the three souls have a relative independence of each other.

By the priority of order, the rational soul should command the sensitive and vegetative souls. However, personal experience shows this not to be the case. It is only with great efforts that the priority of order can be developed. Aquinas makes it clear that man is not an accidental being. If the rational soul were created after the body, then their joining would be accidental. Hence, the idea of the soul is ontologically prior to the body. Aquinas takes pains not to admit the pre-existence of the soul. Nevertheless, we have to conclude that the idea of the soul is prior to the body, so their conjunction is necessary, not accidental.

In any case, the temporal pre-existence of the rational soul in this world makes no sense, which is all that Aquinas means. This is obvious since the rational soul is transcendent to the world process. But this does not preclude that the soul chooses and wills its incarnation. A fortiori, the very nature of the rational soul is intelligence, will, and freedom of choice. To deny this would give credence to the argument that the fetus is not a person since it displays neither intelligence nor will.

The important point is that the rational soul is not an object in the world. Hence, misguided attempts to weigh the soul and take pictures of it leaving the body at death are necessarily measuring something else. As pure consciousness, the rational soul, or Intellect, is inexplicable in natural terms. In his book on the metaphysics of Dante, \textbf{Christian Moevs} writes:

\begin{quotex}
In medieval \textbf{hylomorphism} (the matter-form analysis of reality), pure Intellect (consciousness or awareness) is pure actuality, or form, or Being, or God: it is the self-subsistent principle that spawns or ``contains'' all finite being and experience. Intellect-Being is what is, unqualified, self-subsistent, attributeless, dimensionless. It has no extension in space or time; rather, it projects space-time ``within'' itself, as, analogously, a dreaming intelligence projects a dream-world, or a mind gives being to a thought.
\end{quotex}



\flrightit{Posted on 2012-01-08 by Cologero}

